\section{Conclusion}
Timing channels pose a security threat in modern operating systems. As more of
the simple security issues get patched, more advanced methods of breaking OS
barriers become prevalent. Prior work exists in creating secure microkernels
\cite{ge2019, karlsson2025} as well as promising developments in the user-level
mitigations \cite{cauligi2019}. However, there still exists a gap in how time
protection can be implemented in existing operating systems.

To bridge this gap, this work proposes a domain-level scheduler paired with
flushing during domain switches. Conceptually, this implementation isolates
domains by partitioning execution time into a fixed number of time slices,
which are then dynamically distributed among them. When a domain switch occurs,
the microarchitectural footprint of the previous domain is cleared via the
temporal fence instruction \cite{wistoff2023, wistoff2025}, preventing
cross-domain information leakage through the shared microarchitectural state.

To realize this model, the implementation is built as a modification of the
FreeRTOS RISC-V MPU port, such that MPU protections can be used to completely
separate domains. Because FreeRTOS natively uses a priority-based scheduler, it
is inherently vulnerable to scheduler-based timing attacks where malicious
tasks can infer information from scheduling behaviour. To resolve this, the
domain-level scheduler is introduced on top to manage time slices, while
preserving the native priority-based scheduler within each individual domain.
This approach creates a middle ground, where one can optimize throughput when
time protection is unnecessary. To allow for an intended communication channel
across domain boundaries, the preexisting Queue structure was modified to allow
for cross-domain communication.

Empirical evaluation shows that the implementation is functional. Within the
QEMU-based system emulation, the framework achieves strictly constant-time time
slices, with resulting round-trip times showing a stable cycle count of 50,000,
with a variance of just a single cycle.

However, the implementation is not without its limitations. Notably, it changes
some of the fundamental expectations in FreeRTOS and adds a significant amount
of complexity to the scheduler. There is also known shared state still present,
such as the higher-level caches, which is not isolated through temporal or
spatial partitioning. Some of these areas require more help from hardware,
which would allow for either spatial partitioning or temporal partitioning,
such as in the case of the scratchpad memory used in the S3K kernel
\cite{karlsson2025}. Further, the evaluation itself was significantly limited
by the lack of real-world hardware. The evaluation does not take into account
the underlying microarchitectural state, and as such, the implementation can
only theorize on the efficacy of the temporal fence instruction on domain
switches. It remains as future work to test the implementation on hardware
where such an instruction is present.

Even with this, this thesis has contributed by showing that time protection can
be added into an existing mainstream kernel, while showing intended behaviour
on the QEMU system emulation platform. It shows that time protection is
achievable, but comes with significant trade-offs in complexity, performance
and usability. So while a full domain-level scheduler is most likely not
suitable for more generic operating systems, it remains viable for environments
where a specialized microkernel with security is needed. 

With this in mind, there are still aspects of time protection that could be
used in the more generic operating systems. Specifically, future research can
make use of these concepts:

\begin{itemize}
  \item \textbf{Selective Temporal Partitioning: }Given the low performance
    overhead of the temporal fence instruction (1\% on average), it could be
    used in a generic operating system to create sensitive tasks. The operating
    system could flush the microarchitectural state before and after these
    sensitive tasks are scheduled. This would eliminate their
    microarchitectural footprint while allowing for more freedom in execution
    for the sensitive tasks.
  \item \textbf{Hardware-Assisted partitioning: } In conjunction to this,
    hardware features such as cache partitioning and scratchpad memory could
    guarantee that non-flushable components also remain partitioned from these
    specific sensitive tasks.
  \item \textbf{User level cooperation: } With safe hardware, the user level
    mitigations could cover the remaining attack front in these general purpose
    environments, and could lead to a safe execution environment for known
    sensitive tasks.
\end{itemize}

Finally, to fully guarantee the security contract involved in time protection,
significant work remains in formal verification. As far as this thesis is
concerned, this remains an open challenge.
